\section{Tóm tắt nghiên cứu}

Nghiên cứu này thuộc lĩnh vực \textit{Natural Language Processing} và \textit{Generative AI}, tập trung vào bài toán tối ưu hóa hệ thống \textit{Retrieval-Augmented Generation} (RAG) cho hội thoại hỏi--đáp trên tài liệu kỹ thuật. Trong bài toán hội thoại trên tài liệu, mô hình ngôn ngữ cần trả lời câu hỏi dựa trên các nguồn tài liệu bên ngoài, thay vì chỉ dựa vào kiến thức đã được học trong tham số mô hình. Vì vậy, chất lượng truy xuất tài liệu đóng vai trò quan trọng đối với độ chính xác, tính trung thành với nguồn và khả năng truy vết của câu trả lời.

Các hệ thống RAG hiện nay thường được xây dựng từ nhiều thành phần như phân đoạn văn bản, biến đổi truy vấn, truy xuất dense, truy xuất sparse, truy xuất lai, cấu trúc chỉ mục, sắp xếp lại kết quả và nén ngữ cảnh. Tuy nhiên, nhiều pipeline RAG vẫn được lựa chọn chủ yếu dựa trên kinh nghiệm triển khai, thay vì dựa trên một benchmark có kiểm soát. Điều này tạo ra khoảng trống nghiên cứu: chưa rõ thành phần nào trong pipeline RAG ảnh hưởng mạnh nhất đến chất lượng truy xuất và chất lượng câu trả lời trong bối cảnh hỏi--đáp trên tài liệu kỹ thuật.

Mục tiêu của nghiên cứu là xây dựng một khung đánh giá mang tính kỹ thuật và học thuật nhằm so sánh có hệ thống các cấu hình RAG khác nhau cho bài toán hỏi--đáp dựa trên tài liệu. Khung đánh giá được thiết kế theo năm trục chính: \textit{chunking strategy}, \textit{pre-retrieval method}, \textit{retrieval strategy}, \textit{index structure} và \textit{post-retrieval refinement}. Không gian tổ hợp ban đầu gồm
\[
4 \times 3 \times 3 \times 2 \times 3 = 216
\]
cấu hình. Sau khi loại bỏ các tổ hợp không hợp lệ hoặc dư thừa, nghiên cứu đánh giá 160 pipeline hợp lệ.

Thí nghiệm được thực hiện trên benchmark hỏi--đáp tài liệu, trong đó mỗi câu hỏi có câu trả lời chuẩn và các đoạn tài liệu liên quan. Các chiến lược phân đoạn bao gồm \texttt{fixed\_512}, \texttt{fixed\_1024}, \texttt{recursive} và \texttt{semantic}. Các phương pháp truy xuất bao gồm \texttt{dense\_cosine}, \texttt{sparse\_bm25} và \texttt{hybrid\_rrf}. Thí nghiệm được chia thành hai giai đoạn, trong đó, giai đoạn 1 đánh giá chất lượng truy xuất bằng \textit{Hit Rate@5}, \textit{Recall@5}, \textit{MRR@5} và \textit{NDCG@5}. Giai đoạn 2 chọn các pipeline tốt nhất từ giai đoạn 1 để đánh giá chất lượng sinh câu trả lời bằng các chỉ số gồm \textit{Faithfulness}, \textit{Answer Relevancy} và \textit{Answer Correctness}.

Kết quả thực nghiệm cho thấy \textit{retrieval strategy} là yếu tố có ảnh hưởng lớn nhất đến \textit{Recall@5}, với effect size $\eta^2 = 6.57\%$, cao hơn đáng kể so với \textit{chunking strategy} ($0.15\%$) và \textit{RAG complexity profile} ($0.80\%$). Trong nhóm phương pháp truy xuất, \texttt{sparse\_bm25} đạt mean Recall@5 cao nhất là $0.2655$, tiếp theo là \texttt{hybrid\_rrf} với $0.2559$, trong khi \texttt{dense\_cosine} chỉ đạt $0.0653$. Đối với nhóm chiến lược phân đoạn văn bản, \texttt{recursive} đạt mean Recall@5 cao nhất là $0.1997$, gần với \texttt{fixed\_1024} là $0.1972$. Kết quả cũng cho thấy các pipeline đơn giản, đặc biệt là các cấu hình sử dụng \texttt{fixed\_1024}, \texttt{flat index} và phương pháp truy xuất phù hợp, có thể đạt chất lượng cạnh tranh so với các pipeline phức tạp hơn.

Đóng góp chính của nghiên cứu gồm: (i) xây dựng một khung đánh giá hai giai đoạn cho tối ưu hóa RAG trong bài toán hội thoại trên tài liệu; (ii) đánh giá có hệ thống 160 cấu hình RAG theo năm trục kỹ thuật; (iii) phân tích thống kê bằng ANOVA để xác định mức độ ảnh hưởng của từng nhóm thành phần; và (iv) kết hợp đánh giá ở mức truy xuất và mức sinh câu trả lời nhằm đưa ra kết luận thực nghiệm toàn diện hơn cho việc thiết kế hệ thống RAG.